\documentclass[12pt]{article}
\usepackage[utf8]{inputenc}
\usepackage{amsmath, amssymb, amsfonts}
\usepackage{geometry}
\geometry{a4paper, margin=1in}
\title{MA102 Answer Script Solution}
\author{Roll Number: 210103096}
\date{\today}
\begin{document}
\maketitle
\section*{Solution}

\subsection*{Question 1}

The given matrix is:

\[
A = \left[\begin{array}{ccc}
-1 & 1.5 & 2.5 \\
3.5 & -2 & 0.5 \\
4.5 & 2.5 & -3 \\
-0.5 & -1.5 & 3
\end{array}\right]
\]

To find the basis for the null space of A, we need to perform row reduction on the matrix.

Step-by-step explanation:

1. \textbf{Row Reduction:}
   - Subtract \( \frac{3}{7} \) times row 1 from row 2.
   - Subtract \( \frac{9}{7} \) times row 1 from row 3.
   - Add \( \frac{6}{7} \) times row 1 to row 4.

This results in:

\[
\left[\begin{array}{ccc}
1 & -1.5 & -2.5 \\
0 & 0 & 0 \\
0 & 0 & 0 \\
0 & 0 & 0
\end{array}\right]
\]

2. \textbf{Identify the Basis:}
   - The resulting matrix has only one non-zero row, which means that the null space of A is spanned by the first column vector.

Therefore, the basis for the null space of A is:

\[
\text{Basis B} = \left\{\begin{array}{c}
1 \\
0 \\
0 \\
0
\end{array}\right\}
\]

So, the answer is:

\[
\boxed{\text{Basis B}}
\]

\end{document}